\section{O que foi feito}

Durante a semana, foram realizadas as seguintes atividades técnicas e analíticas:

\subsection{Avaliação comparativa dos datasets candidatos}
Avaliamos os quatro candidatos levantados na fase de pré-projeto com base nos
critérios de independência de registros, combinabilidade matemática de
resultados parciais, relevância contextual e volumetria:
\begin{itemize}
  \item \textbf{Microdados do ENEM (INEP):} Estrutura tabular bem definida
    (formato CSV com delimitador \texttt{;}), contendo entre 3 e 4 milhões de
    participantes por edição anual e volume bruto descompactado de cerca de
    $1{,}5$ a $2{,}0$\,GB. Cada linha representa um inscrito de forma
    totalmente independente, permitindo fatiamento simples por blocos de linhas
    ou bytes.
  \item \textbf{NYC Taxi Trip Records:} Dados com milhões de corridas e
    esquema simples de tarifas. Apesar de excelente para benchmarks, possui
    menor apelo de negócio no contexto acadêmico regional brasileiro.
  \item \textbf{Dumps da Wikipédia:} Modelo clássico de \emph{Word Count}.
    Embora seja didático, o processamento de texto demanda estruturas de dados
    maiores no intercâmbio de mensagens (dicionários extensos de termos), o
    que sobrecarregaria precocemente o canal de comunicação antes dos
    experimentos de heterogeneidade e tolerância a falhas.
  \item \textbf{Dados sintéticos:} Ferramenta essencial de controle de
    volume, viabilizando experimentos estritamente controlados de escalabilidade
    com tamanhos arbitrários ($100$\,MB, $500$\,MB, $1$\,GB, $2$\,GB e
    $4$\,GB) sem variações de latência externa de download.
\end{itemize}

\subsection{Implementação das rotinas de dados}
Criamos o módulo utilitário \texttt{scripts/obter\_dados.py}, responsável por:
\begin{enumerate}
  \item Gerar datasets sintéticos determinísticos com base no esquema de
    colunas do ENEM através do parâmetro \texttt{--seed} (garantindo
    reprodutibilidade estrita entre os integrantes);
  \item Calcular e verificar o hash criptográfico SHA-256 de qualquer arquivo
    de dados (\texttt{--checksum});
  \item Realizar inspeção estatística e contagem rápida de registros
    (\texttt{--inspecionar}), discriminando candidatos presentes e ausentes.
\end{enumerate}

Além do módulo principal, foi criada uma suíte de testes unitários
(\texttt{tests/test\_dados.py}) validando o determinismo de geração de bytes,
a integridade do cabeçalho e as contagens estatísticas de participantes.
